\section{Neural Compute Unit、长期记忆与个性化}

把 LLM 看成 neural compute unit 是后半场最重要的系统类比。模型像计算单元一样接收 token instruction 并输出 token，但它没有自动获得 persistent memory、process isolation、device driver 或 permission model。理解类比的适用范围，才能正确设计 agent architecture。

\subsection{Model as compute unit}

课堂把自然语言 token 与 binary instruction 对比。共同点是输入经过计算变成输出；不同点是 LLM 的行为是概率性的，语义空间开放，输入输出也没有传统 ISA 那样严格的类型与确定性。读图时应把 compute analogy 与 missing systems guarantees 同时保留。

\lecturefigure{slide-090.jpg}{AI model 作为 token-in/token-out neural compute unit}{官方 CS25 V4 slide deck，第 90 页}

\begin{knowledgebox}{读图：CPU analogy 的有效边界}
\begin{tabularx}{\textwidth}{L{0.22\textwidth} L{0.32\textwidth} >{\raggedright\arraybackslash}X}
\toprule
层次 & 类比有效处 & 类比失效处 \\
\midrule
Instruction & Prompt/context 指定本轮计算 & 自然语言含歧义，缺少固定 opcode 与类型系统。 \\
Compute & 参数执行大规模矩阵运算 & 输出是概率分布，不能默认确定性与正确性保证。 \\
Memory & Context 提供工作区 & Context 有长度上限，也不会自动持久化。 \\
I/O & Tool 扩展外部能力 & 权限、异常和副作用必须由系统显式处理。 \\
\bottomrule
\end{tabularx}
\end{knowledgebox}

\teachervoice{讲者把模型称为 neural compute unit，并设想为 Transformer 编写新的 programming language。这个类比的教学价值在于把 prompt、memory、tool 和 variable 变成系统部件，而不是把模型拟人化成完整大脑。}

\subsection{Looped Transformer 与反复计算}

一次 forward pass 对应固定深度计算。Looped Transformer 页提出共享或重复 block，让模型在同一参数上执行更多迭代。它与 reasoning loop 相似：都用更多 computation steps 换取更复杂处理，但必须解决 termination 与 stability。

\lecturefigure{slide-091.jpg}{Looped Transformer 与重复计算模块}{官方 CS25 V4 slide deck，第 91 页}

\begin{importantbox}{Iterative computation}
若同一 block $F_\theta$ 被重复 $K$ 次：
\[
h^{(k+1)}=F_\theta(h^{(k)},x),\qquad k=0,\ldots,K-1.
\]
共享参数降低模型增长速度，却可能带来 optimization、fixed-point stability 和动态 stopping 问题。重复层数多不自动等于 reasoning 更深。
\end{importantbox}

\subsection{Long-term memory 不只是 vector search}

课堂将 persistent memory 类比为 disk：先把文档或经历编码成 embedding 存储，需要时再检索进 context。现有系统常用 nearest-neighbor search，但真实长期记忆还需要时间、层级、关系和更新策略。本节先解释当前 embedding pipeline，再用课堂提出的 open questions 扩展 memory schema。

\lecturefigure{slide-092.jpg}{Long-term memory 的存储、embedding 与检索问题}{官方 CS25 V4 slide deck，第 92 页}

\begin{importantbox}{Memory retrieval score}
对 query embedding $q$ 与 memory item $m_i$，最简相似度可写成
\[
s_i=\frac{q^\top m_i}{\lVert q\rVert_2\lVert m_i\rVert_2},\qquad
\mathcal{I}=\operatorname{TopK}_i(s_i).
\]
这只利用语义相似。真实 memory ranking 还应加入 recency、importance、user scope、causal link、task state 与 access control。
\end{importantbox}

\begin{knowledgebox}{背景概念：Agent memory 的四种层次}
\begin{enumerate}
  \item \textbf{Working memory}：当前 context、scratchpad 和 tool result；
  \item \textbf{Episodic memory}：过去 interaction 与 action trajectory；
  \item \textbf{Semantic memory}：稳定事实、文档和用户知识；
  \item \textbf{Procedural memory}：可复用 workflow、policy 和 skill。
\end{enumerate}
只建一个向量库，通常无法同时满足四种语义。
\end{knowledgebox}

\teachervoice{讲者指出现有 memory 常停留在简单 k-NN。时间一致性、graph structure、online adaptation 和不断变化的数据都没有被自然解决；“接一个 vector database”只是起点。}

\subsection{Personalization：偏好会变化，数据也有边界}

Personalization 页希望 agent 成为用户的数字延伸。它需要学习写作风格、日程、风险偏好和长期目标，同时区分稳定偏好、一次性指令与敏感信息。接下来的两页要同时读 capability 和 data boundary，尤其关注偏好怎样被收集、更新、撤销与覆盖。

\lecturefigure{slide-093.jpg}{User-agent alignment 与 personalization 的目标}{官方 CS25 V4 slide deck，第 93 页}

\lecturefigure{slide-094.jpg}{收集偏好、隐私、变化与 alignment 的挑战}{官方 CS25 V4 slide deck，第 94 页}

\teachervoice{讲者用 hippocampus 和不断变化的个人数据提醒读者：memory 不是静态文档仓库。系统要处理时间顺序、关系结构、在线更新和用户偏好变化。}

\begin{warningbox}{Personalization 的四个常见失败}
\begin{enumerate}
  \item 把一次行为误判为永久偏好；
  \item 将不同用户或 workspace 的 memory 混用；
  \item 为了“更懂用户”无限保存敏感数据；
  \item 用户偏好与安全、法律或组织 policy 冲突时盲目服从。
\end{enumerate}
系统需要 consent、scope、retention、edit/delete 与 policy precedence。
\end{warningbox}

\begin{importantbox}{Preference update 需要时间衰减与确认}
可把 preference $p$ 的 score 写成
\[
S(p)=\sum_j w_j\,e^{-\lambda\Delta t_j}\,\operatorname{evidence}_j(p),
\]
其中 $w_j$ 表示证据可靠性，$\Delta t_j$ 是时间间隔，$\lambda$ 控制衰减。高风险偏好不应只靠隐式行为推断，而应要求显式确认。
\end{importantbox}

\subsection{本章小结}

Neural-compute analogy 让 model、context、memory 与 tool 的边界更清楚，但不能把 LLM 当成具有确定 ISA 和隔离保证的 CPU。长期记忆需要层级、时间和权限，个性化需要 consent、更新和删除机制。系统能力越贴近用户，memory governance 越成为核心设计。

